\newpage
\phantomsection
\label{prf:telescoping-identity}
\LRAProofFor{thm:telescoping-identity}

\begin{remark*}[Return]
\hyperref[thm:telescoping-identity]{Return to Theorem}
\end{remark*}

\begin{theorem*}[Telescoping Identity]
For any function \(f : \mathbb{Z} \to \mathbb{R}\) and integers \(a \le b\),
\[
\sum_{k=a}^{b-1} (f(k+1) - f(k)) = f(b) - f(a).
\]
\end{theorem*}

\begin{proof}[Professional Standard Proof]
\LRAProofBodyStart
Expanding the sum gives
\[
\begin{aligned}
(f(a+1)-f(a)) &+ (f(a+2)-f(a+1)) \\
              &+ \cdots + (f(b)-f(b-1)).
\end{aligned}
\]
Every intermediate term \(f(a+1), f(a+2), \dots, f(b-1)\) appears once
with a positive sign and once with a negative sign, so they all cancel.
The remaining terms are \(f(b)\) and \(-f(a)\), giving \(f(b)-f(a)\).
\end{proof}

\begin{proof}[Detailed Learning Proof]
\LRAProofBodyStart
Expanding the sum gives
\[
\begin{aligned}
(f(a+1)-f(a)) &+ (f(a+2)-f(a+1)) \\
              &+ \cdots + (f(b)-f(b-1)).
\end{aligned}
\]
Every intermediate term \(f(a+1), f(a+2), \dots, f(b-1)\) appears once
with a positive sign and once with a negative sign, so they all cancel.
The remaining terms are \(f(b)\) and \(-f(a)\), giving \(f(b)-f(a)\).
\end{proof}

\begin{remark*}[Proof structure]
Expand the finite sum and cancel adjacent endpoint terms.
\end{remark*}

\begin{dependencies}
\begin{itemize}
  \item \hyperref[def:sigma-summation-notation]{Summation Notation}
\end{itemize}
\end{dependencies}
\clearpage
